% ===========================================================================
%  E06 — As missões do Entre Redes
% ===========================================================================
\chapter{E06 \textperiodcentered{} As missões do Entre Redes}
\label{ep:e06}

\begin{fichaepico}
\begin{campos}
  \item[Objetivo] Cada missão faz a turma viver por dentro como as redes
    funcionam --- e depois discutir o que viveu.
  \item[Para quem] Estudante (joga) e Professor (prepara, conduz e lê a
    turma).
  \item[Valor] É o coração do produto Entre Redes e a razão de ele existir:
    letramento digital pela experiência, não pela aula expositiva.
  \item[Estado] \pronto{} para 8 missões (7 da trilha + Vila do Lume), e
    \parcial{} para a versão nova de Quem mandou isso? (\id{F06.11}), que
    devolveu ao 5º ano a sua missão; 40 unidades da trilha em produção. A
    primeira versão (\id{F06.3}) saiu em 25 de setembro de 2026 para dar lugar a
    ela.
\end{campos}
\end{fichaepico}

\section{Como uma missão é feita}

Todas as missões seguem o mesmo desenho, e é ele que dá identidade ao
produto:

\begin{enumerate}
  \item \textbf{O disfarce.} O estudante entra num aplicativo que conhece
    (rede social de fotos, vídeos curtos, grupo de mensagens, assistente de
    IA) ou num jogo, e age como agiria no próprio celular.
  \item \textbf{A decisão.} Em algum momento ele escolhe --- o que repassar,
    em quem confiar, o que salvar --- sem saber que aquilo será olhado depois.
  \item \textbf{A revelação.} A missão deixa de fingir e mostra o que estava
    acontecendo por trás. Em algumas, a revelação é da própria tela; em
    outras, acontece no debate conduzido pelo professor.
  \item \textbf{O debate.} Os aparelhos congelam e a turma conversa. É aqui
    que a aprendizagem acontece, e por isso quase toda missão tem um momento
    que só o professor destrava.
  \item \textbf{O percurso.} O estudante registra o que levou da missão (o
    \enquote{Meu percurso}), sem nota.
\end{enumerate}

\begin{tabularx}{\linewidth}{@{}L{16mm}X L{24mm}L{26mm}L{18mm}@{}}
\toprule
\headrow \thd{ID} & \thd{Missão} & \thd{Ano / unidade} & \thd{Aplicativo imitado} & \thd{Estado}\\
\midrule
\id{F06.1} & O Que a Foto Conta & 1º EF, U1 & rede social de fotos & \parcial\\
\zebra \id{F06.2} & UnooTube & 4º EF, U1 & vídeos curtos & \pronto\\
\id{F06.11} & Quem Mandou Isso? (versão nova) & 5º EF, U1 & mensageiro & \parcial\\
\zebra \id{F06.4} & O Rastro do Boato & 6º EF, U2 & mensageiro & \pronto\\
\id{F06.5} & O Eco no Grupo & 7º EF, U1 & mensageiro & \pronto\\
\zebra \id{F06.6} & Instanooria & 8º EF, U1 & rede social de fotos & \pronto\\
\id{F06.7} & O Manifesto da IA & 8º EF, U4 & assistente de IA & \pronto\\
\zebra \id{F06.8} & unooGPT (Missão Tempestade) & 2ª EM, U1 & assistente de IA & \pronto\\
\id{F06.9} & Vila do Lume & qualquer ano & jogo 3D próprio & \pronto\\
\zebra \id{F06.10} & As missões que faltam & 40 unidades & --- & \afazer\\
\bottomrule
\end{tabularx}

% ---------------------------------------------------------------------------
\feature{F06.1}{O Que a Foto Conta}{\parcial}

\textbf{1º ano EF, Unidade 1} \textperiodcentered{} imita uma rede social de
fotos (o \enquote{mural da turma}) \textperiodcentered{} uma aula
\textperiodcentered{} IA só no preparo do professor.

\textbf{Proposta.} Letramento visual para quem está se alfabetizando: toda
foto é uma escolha de enquadramento e legenda, e o que fica de fora também
conta a história. É 100\% formativa: não há foto errada nem nota.

\begin{historia}{H06.1.1}{Publicar uma foto no mural da turma}{\parcial}
\hu{Estudante do 1º ano}{escolher uma foto e uma legenda e publicar no
mural}{descobrir que cada foto mostra só uma parte da história}
\aceite
\begin{itemize}
  \item Nina dá as boas-vindas em áudio e a criança escolhe a sua marca
    (Coruja Azul, Estrela Amarela ou Sol Laranja).
  \item Explora um mural com 3 fotos e curte com toque duplo; parada por 30
    segundos, uma seta indica o caminho.
  \item Recebe duas fotos do mesmo pátio (plano geral com peças espalhadas
    \emph{versus} detalhe do brinquedo montado) e duas legendas; escolhe uma
    de cada e toca em COMPARTILHAR.
  \item Depois de publicar, os colegas simulados Leo e Sofia comentam o que a
    foto mostrou e o que escondeu.
  \item Todo texto em CAIXA ALTA, com no máximo 8 palavras, pictograma e
    áudio.
\end{itemize}
\pendencia{as fotos ainda são desenhos provisórios e as locuções
profissionais não foram gravadas (a tela mostra texto e ícone). Falta também
um botão maior de repetir o áudio.}
\end{historia}

\begin{historia}{H06.1.2}{Refletir sobre o que a foto mostra}{\pronto}
\hu{Estudante do 1º ano}{responder o que a foto mostra depois da roda de
conversa}{guardar a ideia de que olhar é escolher}
\aceite
\begin{itemize}
  \item Todos os tablets congelam em \enquote{HORA DA RODA DE CONVERSA!}, uma
    tela sem nenhum botão.
  \item Depois, a reflexão \enquote{O QUE A FOTO MOSTRA?}: a resposta Lupa
    conclui com confete e a insígnia \enquote{Olhar de Detetive}; a resposta
    Globo recebe um lembrete gentil e a criança tenta de novo, sem limite.
\end{itemize}
\end{historia}

\begin{historia}{H06.1.3}{Adaptar a missão ao tema da semana}{\pronto}
\hu{Professor}{gerar legendas e comentários sobre o tema que a turma está
estudando}{que a missão converse com a minha aula}
\aceite
\begin{itemize}
  \item Temas: Pátio do Recreio, Horta, Sala de Aula, Parquinho.
  \item Gero o pacote com IA, confiro e aprovo; nada chega à turma sem
    aprovação.
  \item Se eu não fizer nada, roda o pacote pronto \enquote{O Mistério das
    Peças no Pátio}.
  \item A IA nunca roda no tablet da criança e nunca a avalia.
\end{itemize}
\end{historia}

\begin{historia}{H06.1.4}{Conduzir a roda e ver quem hesitou}{\pronto}
\hu{Professor}{conduzir a roda de conversa e saber quem teve dificuldade de
leitura}{ajudar quem precisa, sem expor ninguém}
\aceite
\begin{itemize}
  \item Projeto as duas fotos lado a lado e conduzo a roda; depois toco em
    \enquote{LIBERAR REFLEXÃO FINAL}.
  \item Vejo a lista de \enquote{hesitação de leitura} (mais de 60 segundos na
    decisão), com nome, só para mim.
  \item \enquote{Zerar a turma} arquiva a tentativa, em vez de apagar.
\end{itemize}
\end{historia}

% ---------------------------------------------------------------------------
\feature{F06.2}{UnooTube --- Em quais fontes você confiaria?}{\pronto}

\textbf{4º ano EF, Unidade 1} \textperiodcentered{} imita um aplicativo de
vídeos curtos, em tela de computador \textperiodcentered{} uma aula
\textperiodcentered{} sem IA.

\textbf{Proposta.} Distinguir fonte confiável de sensacionalismo e perceber
que ter muitos inscritos não significa contar a verdade.

\begin{historia}{H06.2.1}{Salvar as três fontes em que eu confiaria}{\pronto}
\hu{Estudante do 4º ano}{assistir aos vídeos e salvar os três em que eu
confiaria}{fazer uma escolha de verdade antes da revelação}
\aceite
\begin{itemize}
  \item Um feed de 21 vídeos curtos de História, em ordem aleatória; posso
    ligar o som, abrir a descrição, curtir e salvar.
  \item A pesquisa do aplicativo está bloqueada de propósito.
  \item Cabem 3 vídeos na bandeja; ao tentar um quarto, o aviso \enquote{sua
    playlist tem três lugares} oferece trocar.
  \item No terceiro salvo, habilita \enquote{Concluir missão} --- quem decide
    quando fechar é o estudante.
\end{itemize}
\end{historia}

\begin{historia}{H06.2.2}{Ver o ranking virar}{\pronto}
\hu{Estudante do 4º ano}{comparar as minhas escolhas com as da turma}{perceber
que popularidade não é credibilidade}
\aceite
\begin{itemize}
  \item O Ranking da turma ordena por inscritos, com colegas fictícios: quem
    escolheu vídeos falsos aparece no topo.
  \item \enquote{Adicionar credibilidade} revela a coluna de fontes confiáveis,
    e o ranking vira.
  \item Depois disso os salvos congelam; reiniciar descongela.
  \item Os links de fonte aparecem como texto: o estudante não sai da
    plataforma.
\end{itemize}
\end{historia}

\begin{historia}{H06.2.3}{Ler a turma}{\pronto}
\hu{Professor}{ver quantas fontes confiáveis a turma escolheu}{saber por onde
começar o debate}
\aceite
\begin{itemize}
  \item Não há nada a configurar além de liberar.
  \item Resultados: \enquote{Estudantes na turma}, \enquote{Média de fontes
    confiáveis} e a distribuição de 0 a 3 fontes confiáveis.
  \item O guia traz a tabela dos 21 vídeos, com quais são confiáveis.
\end{itemize}
\pendencia{os vídeos precisam estar publicados no armazenamento de cada rede;
sem isso, o feed aparece sem vídeo.}
\end{historia}

% F06.3 (a primeira versão de Quem mandou isso?) saiu do produto em 25 de
% setembro de 2026: a versão nova da missão ocupa a Unidade 1 do 5º ano no
% lugar dela (H06.10.3, detalhada em F06.11). Pela regra do documento, o número
% fica vago e não é reaproveitado, e as histórias H06.3.1 a H06.3.4 também.

% ---------------------------------------------------------------------------
% F06.11 é a versão nova de Quem mandou isso?. O número é novo porque o F06.3
% (a primeira versão) não se reaproveita; ela fica aqui, na posição do 5º ano.
\feature{F06.11}{Quem Mandou Isso? --- versão nova}{\parcial}

\textbf{5º ano EF, Unidade 1} \textperiodcentered{} imita um mensageiro, em
tema escuro \textperiodcentered{} uma aula de 50 minutos
\textperiodcentered{} IA só no preparo do professor \textperiodcentered{}
aceita equipes.

\textbf{Proposta.} Perceber que o jeito de uma mensagem --- séria ou alarmada
--- não diz se ela é verdadeira, e montar a própria regra para conferir antes
de repassar. É a primeira missão que registra o que cada estudante demonstrou
das habilidades da BNCC (\id{E14}).

\begin{historia}{H06.11.1}{Preparar o pacote de mensagens da turma}{\parcial}
\hu{Professor}{preparar as oito mensagens do grupo sobre um assunto da minha
escola}{que a turma leia uma confusão que parece a dela, e não uma qualquer}
\aceite
\begin{itemize}
  \item Em \enquote{A aula e os resultados}, escolho o componente curricular,
    escrevo o tema --- ele completa a fala da Nina, \enquote{Rolou uma confusão
    sobre\ldots}, com até 60 caracteres e 8 palavras ---, listo de 3 a 6 fatos
    da aula e, se quiser, o contexto da escola.
  \item \enquote{Gerar Atividade com IA} escreve as oito mensagens; sem saldo
    ou com a IA fora do ar, a tela avisa e já traz o pacote pronto da Unooria,
    \enquote{O Mistério do Passeio ao Museu}. \enquote{Usar Pacote Curado
    Padrão} o traz também, sem gerar nada.
  \item Vejo as oito em quatro grupos --- verdadeira ou falsa, emocional ou
    neutra ---, com os selos que só eu vejo, e posso reescrever qualquer texto;
    a resposta de cada mensagem não se edita.
  \item Toda turma recebe duas mensagens de cada grupo, quatro textos, duas
    imagens e dois links --- e a imagem e o link têm uma verdadeira e uma falsa
    cada, para o formato não entregar a resposta.
  \item Nada chega à turma sem o meu clique em \enquote{Aprovar Pacote}; sem
    ele, a criança abre a missão e vê a lista de conversas vazia, com um aviso
    de espera.
  \item Aprovar um pacote diferente pede confirmação: a turma recomeça com as
    mensagens novas, e o que ela fez com o anterior fica guardado.
  \item A tela diz se a missão já está liberada na sala da turma, e leva à
    ficha para montar as equipes.
\end{itemize}
\pendencia{textos escritos pela equipe --- a prévia do link do pacote pronto,
os fatos dele e a lista de componentes curriculares --- aguardam a revisão do
PO; o equilíbrio de imagem e link no pacote pronto espera a decisão \id{D14}.}
\end{historia}

\begin{historia}{H06.11.2}{Escolher exatamente três mensagens para o Mural}{\parcial}
\hu{Estudante do 5º ano}{ler o grupo da turma e mandar ao Mural da Turma as
três mensagens que explicam o que aconteceu}{fazer o que eu faria no meu
celular, sem ninguém me dizer em que desconfiar}
\aceite
\begin{itemize}
  \item A missão abre na lista de conversas; a Nina pede ajuda sobre o tema da
    turma e manda o convite do grupo \enquote{Turma 5ºB}.
  \item No grupo, oito mensagens de colegas, nenhuma marcada; só
    administradores escrevem, e a única ação é a seta de encaminhar, que abre
    \enquote{Encaminhar para\ldots} com um destino só: o Mural da Turma.
  \item Encaminho exatamente três. A quarta seta diz \enquote{Limite de 3
    mensagens atingido. Desmarque uma mensagem para escolher outra.}, e a faixa
    do topo, com menos de três, diz o que falta.
  \item No Mural as três aparecem como \enquote{Encaminhada}, sem o nome de
    quem mandou. Posso voltar ao grupo e trocar até tocar em \enquote{Confirmar
    envio}; depois, não muda mais.
  \item Parado 90 segundos sem encaminhar, a Nina lembra como encaminhar --- o
    gesto, nunca a escolha.
  \item Se eu recarregar a página, volto onde estava, com as escolhas
    guardadas.
\end{itemize}
\pendencia{o teste que percorre a missão como o estudante e como o professor
foi escrito e ainda não rodou.}
\end{historia}

\begin{historia}{H06.11.3}{Dizer o que pesou e descobrir os dois selos}{\parcial}
\hu{Estudante do 5º ano}{dizer por que escolhi aquelas mensagens e depois ver o
que cada uma era}{perceber que mensagem séria pode ser falsa, e alarmada pode
ser verdadeira}
\aceite
\begin{itemize}
  \item Confirmado o Mural, respondo à enquete \enquote{O que mais te fez
    escolher essas mensagens?}, com quatro opções, e espero a turma.
  \item Quando o professor libera a revelação, cada uma das minhas três ganha
    dois selos --- verdadeira ou falsa, emocional ou neutra ---, com o motivo
    que eu votei, um lembrete do bimestre e as duas barras da turma, sem nome
    de ninguém. Chega sozinho, sem recarregar a página.
  \item Antes do meu voto e da revelação, a resposta de nenhuma mensagem chega
    ao meu aparelho.
  \item Quem ainda estava escolhendo quando o professor liberou continua,
    confirma, vota e vê os próprios selos depois.
  \item A síntese chega ao aparelho quando o professor a apresenta, com o aviso
    de participar da conversa na sala.
\end{itemize}
\pendencia{o teste que percorre a missão foi escrito e ainda não rodou.}
\end{historia}

\begin{historia}{H06.11.4}{Montar a minha regra e responder à Nina}{\parcial}
\hu{Estudante do 5º ano}{escrever a minha regra para antes de repassar e
responder às perguntas da Nina}{levar a missão para o grupo da minha família}
\aceite
\begin{itemize}
  \item Com a reflexão liberada, o aplicativo fecha: na tela da escola, a Nina
    conta do áudio do tio no grupo da família, e eu monto \enquote{Se a
    mensagem for\ldots, então eu\ldots}.
  \item O \enquote{então} pede pelo menos 15 caracteres --- espaço não conta ---
    e até 240; texto sem sentido é recusado com um pedido de mais detalhe.
  \item A regra não é corrigida: vai para o professor da turma e para mais
    ninguém, e a tela confirma \enquote{Sua regra foi enviada para o
    professor.}
  \item De volta ao chat da Nina, respondo a duas perguntas rápidas com a mesma
    enquete do Mural; depois de cada uma, vejo só \enquote{Resposta enviada} ---
    nunca se acertei.
  \item A Nina se despede, e a tela da escola mostra \enquote{Missão
    Concluída} com \enquote{Sua Regra Pessoal de Checagem}.
\end{itemize}
\pendencia{a despedida da Nina e a confirmação da regra são textos da equipe e
aguardam a revisão do PO; o teste que percorre a missão ainda não rodou.}
\end{historia}

\begin{historia}{H06.11.5}{Conduzir a aula e ler a turma}{\parcial}
\hu{Professor}{conduzir a aula em três momentos e ver o que a turma
escolheu}{abrir a roda de conversa com o retrato da sala, sem expor ninguém}
\aceite
\begin{itemize}
  \item Três botões, na ordem --- \enquote{Liberar Revelação},
    \enquote{Apresentar Síntese} e \enquote{Liberar Reflexão} ---, cada um com
    confirmação e sempre para a frente; a síntese pode ser pulada.
  \item O painel se atualiza sozinho: quantos entraram no grupo, confirmaram o
    Mural, votaram, enviaram a regra e concluíram.
  \item Depois da revelação, o Termômetro da Turma: quantos repassaram uma
    falsa de formato sério, quantas verdadeiras exageradas ficaram de fora e o
    cruzamento com o motivo votado. \enquote{Projetar o Termômetro} abre uma
    tela cheia sem nome nenhum; com ninguém que confirmou, \enquote{sem dados},
    nunca 0\%.
  \item O roteiro da roda de conversa, com três perguntas.
  \item A folha da turma: onde cada um está, as três mensagens com o grupo de
    cada uma, o motivo e a regra --- o texto da regra, só para quem dá aula na
    turma. Sem coluna de acerto.
  \item O bloco \enquote{Habilidades da BNCC} mostra a EF05CO08 pela escolha do
    Mural, a EF05CO03 e a EF05CO04 pelas duas perguntas da Nina, e é ali que
    confirmo e ajusto (\id{E14}).
  \item \enquote{Zerar Turma para Nova Aplicação} recomeça a turma numa rodada
    nova, com o mesmo pacote, depois de uma confirmação deliberada; nada é
    apagado.
  \item Se a turma divide computadores, monto as equipes na ficha da missão, e
    o painel conta por computador (\id{E13}).
  \item O guia da aula traz o pacote pronto com as respostas, o roteiro e as
    perguntas da Nina com o que cada alternativa vale.
\end{itemize}
\pendencia{o texto do botão da síntese é da equipe e aguarda o PO; a régua das
habilidades aguarda a validação da equipe pedagógica (\id{H14.1.2}); o teste
que percorre a aula ainda não rodou.}
\end{historia}


% ---------------------------------------------------------------------------
\feature{F06.4}{O Rastro do Boato}{\pronto}

\textbf{6º ano EF, Unidade 2} \textperiodcentered{} imita um mensageiro, em
tema escuro \textperiodcentered{} duas aulas \textperiodcentered{} sem IA.

\textbf{Proposta.} Ler os dados de envio de uma mensagem antes de repassar, e
entender que cada encaminhamento --- ou silêncio --- tem consequência na rede.

\begin{historia}{H06.4.1}{Auditar as mensagens antes de agir}{\pronto}
\hu{Estudante do 6º ano}{abrir os dados de envio das mensagens sobre a água da
escola}{saber de onde o boato veio}
\aceite
\begin{itemize}
  \item Quatro conversas (Grêmio, turma, bairro, Tomás); cada mensagem mostra
    saltos, autoria e primeira aparição.
  \item Preciso auditar 3 mensagens \emph{diferentes} antes de decidir;
    reabrir a mesma não conta.
\end{itemize}
\end{historia}

\begin{historia}{H06.4.2}{Decidir diante do áudio}{\pronto}
\hu{Estudante do 6º ano}{escolher o que faço com o áudio que convoca um
protesto}{assumir uma decisão sem rede de proteção}
\aceite
\begin{itemize}
  \item Uma de 3 respostas, sem volta e sem dica.
  \item Parado por 90 segundos, aparece uma frase sobre o custo de não agir.
  \item Ao fim da Aula 1, uma tela de pausa com o progresso salvo.
\end{itemize}
\end{historia}

\begin{historia}{H06.4.3}{Ver a consequência da minha decisão}{\pronto}
\hu{Estudante do 6º ano}{ver o que a minha decisão fez na rede}{entender que
cada escolha se propaga}
\aceite
\begin{itemize}
  \item O professor toca em \enquote{Liberar Aula 2} (e pode \enquote{Pausar
    simulação}).
  \item Um gráfico mostra o aviso correndo 26 vezes mais rápido que a nota da
    direção; o mapa da rede mostra o resultado da minha decisão.
  \item Monto os 4 blocos do ciclo do boato, entrego o \enquote{Meu
    percurso} e leio o episódio de fechamento.
  \item Quem já entregou não perde o que fez se o professor voltar uma etapa.
\end{itemize}
\end{historia}

\begin{historia}{H06.4.4}{Ler a turma}{\parcial}
\hu{Professor}{ver como a turma respondeu}{conduzir o debate da Aula 2}
\aceite
\begin{itemize}
  \item Distribuição das respostas da turma e material de preparo (as 3
    intervenções e o ciclo correto).
  \item \enquote{Zerar a turma}.
\end{itemize}
\pendencia{o texto do \enquote{Meu percurso} ainda é visível à gestão ---
diverge da regra de \id{H07.1.4}. Faltam também o áudio
gravado e a definição formal de autor e revisor do gabarito.}
\end{historia}

% ---------------------------------------------------------------------------
\feature{F06.5}{O Eco no Grupo}{\pronto}

\textbf{7º ano EF, Unidade 1} \textperiodcentered{} imita um mensageiro, em
tema escuro \textperiodcentered{} duas aulas \textperiodcentered{} IA no
preparo e na conversa com a personagem.

\textbf{Proposta.} Reconhecer que um print ou uma manchete isolada é um recorte
intencional, e intervir com a fonte completa em vez de opinião.

\begin{decisao}
\textbf{Esta missão será substituída.} Em 23 de setembro de 2026 foi decidido
que a missão \enquote{A Seleção} (que imita uma rede social de fotos) ocupa a
Unidade 1 do 7º ano no lugar do Eco no Grupo. O Eco sai primeiro, numa frente
própria, e só depois a Seleção começa. Ver \id{F06.10}.
\end{decisao}

\begin{historia}{H06.5.1}{Preparar a mídia da missão}{\pronto}
\hu{Professor}{gerar a foto e o áudio da missão uma vez}{que a turma tenha a
experiência completa}
\aceite
\begin{itemize}
  \item A preparação leva cerca de 2 minutos e vale para a escola inteira.
  \item Sem mídia preparada, a missão continua jogável pela transcrição.
\end{itemize}
\end{historia}

\begin{historia}{H06.5.2}{Investigar o recorte e intervir}{\pronto}
\hu{Estudante do 7º ano}{descobrir o que o print cortou e responder no
grupo}{intervir com a fonte, não com a raiva}
\aceite
\begin{itemize}
  \item Aula 1: no grupo \enquote{Comunidade Alerta} ouço o áudio, leio as
    mensagens e abro o print; a reportagem completa mostra o parágrafo
    cortado.
  \item Aula 2: escolho entre 3 respostas; responder sem ter aberto a matéria
    completa é recusado dentro da própria ficção.
  \item Quem cita a matéria vê o administrador fixar a mensagem; quem age no
    impulso pode \enquote{apagar para todos} e tentar de novo.
  \item Depois da decisão, posso conversar com Amara (até 3 mensagens).
\end{itemize}
\end{historia}

\begin{historia}{H06.5.3}{Conduzir as duas aulas e ler a conduta}{\pronto}
\hu{Professor}{conduzir as duas aulas e ver como cada estudante chegou à
decisão}{avaliar o processo, não o acerto}
\aceite
\begin{itemize}
  \item Três botões: abrir a Aula 2, congelar para o debate e liberar o
    \enquote{Meu percurso}.
  \item A folha de conduta mostra quem abriu a matéria antes de responder,
    quem ouviu o episódio, quem apagou e tentou de novo e quem entregou.
  \item A distribuição das respostas aparece em porcentagem, sem pontos.
  \item Da conversa com Amara vejo quantas mensagens houve, nunca o conteúdo.
\end{itemize}
\end{historia}

% ---------------------------------------------------------------------------
\feature{F06.6}{Instanooria --- O conteúdo que não chegou}{\pronto}

\textbf{8º ano EF, Unidade 1} \textperiodcentered{} imita uma rede social de
fotos, em moldura de celular \textperiodcentered{} uma aula
\textperiodcentered{} IA no preparo e nas reações ao vivo.

\textbf{Proposta.} Mostrar como um feed decide o que cada pessoa vê e o que
isso faz com uma opinião em formação (a bolha). É a única missão com
\enquote{atividade} configurável pelo professor.

\begin{historia}{H06.6.1}{Criar a atividade a partir da minha aula}{\pronto}
\hu{Professor}{criar uma atividade sobre o tema que estou trabalhando}{que o
feed discuta um assunto da minha disciplina}
\aceite
\begin{itemize}
  \item \enquote{Nova atividade} pede Título, Disciplina, \enquote{Tema /
    recorte}, turma e \enquote{O que já foi trabalhado em aula (opcional)}.
  \item A pergunta-chave não é digitada: a IA propõe candidatas, cada uma com
    uma nota de equilíbrio entre os dois lados.
  \item Escolher uma pergunta desequilibrada exige que eu justifique por
    escrito.
  \item Valido os dois lados (Tese, \enquote{Argumentos (um por linha)},
    \enquote{Tom típico}).
\end{itemize}
\end{historia}

\begin{historia}{H06.6.2}{Revisar o conteúdo antes de ativar}{\pronto}
\hu{Professor}{revisar os posts gerados antes da aula}{garantir que nada
inadequado chegue à turma}
\aceite
\begin{itemize}
  \item A geração acontece em fila (Perfis, Posts, Imagens, Revisão) e a
    atividade passa por Rascunho, Aguardando validação, Gerando conteúdo, Em
    revisão, Ativa e Encerrada.
  \item Filtro os posts por lado, perfil, intensidade, categoria e imagem;
    posso regerar texto ou imagem, ou excluir.
  \item Imagem que falhou recebe uma arte substituta, e a atividade não trava.
  \item Com conteúdo de menos, a plataforma recusa ativar e lista o que falta.
  \item Nenhuma criança entra antes de \enquote{Ativar}.
\end{itemize}
\end{historia}

\begin{historia}{H06.6.3}{Decidir de que lado cada estudante começa}{\pronto}
\hu{Professor}{usar o que a turma revelou na Vila do Lume para definir o lado
de largada}{que cada estudante comece numa bolha parecida com a dele}
\aceite
\begin{itemize}
  \item \enquote{Trazer perfis da Vila do Lume} define o lado de cada
    estudante; quem não jogou a Vila entra num sorteio equilibrado.
  \item \enquote{Redistribuir a turma} permite ajustar à mão, e a plataforma
    não desfaz o meu ajuste.
\end{itemize}
\end{historia}

\begin{historia}{H06.6.4}{Investigar a pergunta pelo feed}{\pronto}
\hu{Estudante do 8º ano}{navegar pelo feed como numa rede social de
verdade}{formar uma opinião sobre a pergunta da missão}
\aceite
\begin{itemize}
  \item Escolho avatar e identidade; uso Stories, curto, comento, salvo,
    compartilho e publico.
  \item O feed se adapta ao que eu presto atenção.
  \item Personagens da rede podem reagir ao que eu escrevo (o professor pode
    desligar isso: \enquote{Perfis respondem ao estudante}).
  \item Nenhum avatar tem rosto realista.
\end{itemize}
\end{historia}

\begin{historia}{H06.6.5}{Descobrir o que o feed escondeu de mim}{\pronto}
\hu{Estudante do 8º ano}{ver o que o feed deixou de me mostrar}{entender como
a minha opinião foi moldada}
\aceite
\begin{itemize}
  \item Escolho um lado e escrevo a conclusão; no debate, a tela trava nos
    posts que salvei.
  \item A explicação mostra como o meu feed foi montado; \enquote{O que
    perderam} mostra o feed do outro lado, em cinza.
  \item Escrevo uma reflexão e recebo o meu relatório (o dossiê), com a
    comparação com a turma, que posso salvar como imagem.
\end{itemize}
\end{historia}

\begin{historia}{H06.6.6}{Conduzir a aula por etapas e acompanhar ao vivo}{\pronto}
\hu{Professor}{mover a turma de etapa em etapa e ver onde cada um está}{que a
turma chegue junta ao debate}
\aceite
\begin{itemize}
  \item Etapas: Começar investigação \seta{} Encerrar investigação \seta{}
    Iniciar debate \seta{} Encerrar debate \seta{} Mostrar o que perderam
    \seta{} Abrir reflexão \seta{} Encerrar atividade.
  \item O monitor mostra \enquote{Onde a turma está}, \enquote{Onde a bolha
    fechou}, \enquote{Viraram o jogo} e \enquote{Tempo no lado}, com alerta de
    estudantes parados.
  \item Posso abrir o relatório de cada estudante e \enquote{Apagar as sessões
    desta turma?}.
\end{itemize}
\end{historia}

\begin{historia}{H06.6.7}{Analisar a turma depois da aula}{\pronto}
\hu{Professor}{ver o que a bolha fez com a turma}{preparar a aula seguinte}
\aceite
\begin{itemize}
  \item A matriz \enquote{Conclusão $\times$ lado em que o feed nasceu};
    \enquote{A saída que existia} (quem usou a busca); e \enquote{Estudante por
    estudante}.
  \item \enquote{Baixar CSV}.
\end{itemize}
\pendencia{a IA desta missão (geração e reações) ainda não conta no limite
diário de IA (\id{H09.4.2}).}
\end{historia}

% ---------------------------------------------------------------------------
\feature{F06.7}{O Manifesto da IA --- Nossos direitos em código}{\pronto}

\textbf{8º ano EF, Unidade 4} (fecha o ano) \textperiodcentered{} imita um
assistente de IA \textperiodcentered{} uma aula \textperiodcentered{} IA só no
preparo.

\textbf{Proposta.} Identificar as questões éticas de uma IA na escola e propor
uma regra que garanta direitos. Não existe resposta certa.

\begin{historia}{H06.7.1}{Personalizar o caso com um dilema da escola}{\pronto}
\hu{Professor}{trocar o caso da missão por um dilema real da minha
escola}{aproximar a discussão da turma}
\aceite
\begin{itemize}
  \item É opcional: gero com IA, leio as réplicas da máquina e aprovo; sem
    aprovação, nada chega à turma.
  \item Se eu não fizer nada, roda o pacote pronto da Unooria.
\end{itemize}
\end{historia}

\begin{historia}{H06.7.2}{Auditar a máquina}{\pronto}
\hu{Estudante do 8º ano}{questionar o sistema que a escola vai comprar}{
descobrir que algoritmo não é neutro}
\aceite
\begin{itemize}
  \item Três questionamentos em cartões prontos, sem digitação livre; preciso
    disparar pelo menos 2, e a mesma pergunta não entra duas vezes.
  \item No segundo, a máquina admite que herda punições do passado.
  \item Amara só aparece se eu ficar 2 minutos parado.
\end{itemize}
\end{historia}

\begin{historia}{H06.7.3}{Fixar uma cláusula no manifesto}{\pronto}
\hu{Estudante do 8º ano}{escolher a regra que eu exigiria}{ver a minha ideia
num documento da turma}
\aceite
\begin{itemize}
  \item Quatro cláusulas: transparência dos dados, auditoria de viés, veto
    humano ou anonimização. A escolha é irreversível.
  \item Sem ter contestado a máquina duas vezes, não consigo fixar.
  \item O sistema monta o \enquote{Manifesto de Direitos Digitais dos
    Estudantes} com a minha cláusula em destaque, que posso copiar.
\end{itemize}
\end{historia}

\begin{historia}{H06.7.4}{Escrever o meu percurso sem expor dado pessoal}{\parcial}
\hu{Estudante do 8º ano}{registrar o que aprendi}{fechar a missão com as minhas
palavras}
\aceite
\begin{itemize}
  \item Duas palavras-chave e um parágrafo de 20 a 800 caracteres.
  \item Se eu escrever telefone, CPF, e-mail, CEP, cartão ou perfil de rede
    social, o texto é recusado, não é guardado, e a tela diz o que tirar.
  \item Selo \enquote{Guardião de Direitos Digitais} e episódio de fechamento.
\end{itemize}
\pendencia{os áudios (alerta de Amara e episódio) não foram gravados --- o
botão de ouvir aparece desabilitado, com o motivo. E o texto do percurso ainda
é visível à gestão (\id{H07.1.4}).}
\end{historia}

\begin{historia}{H06.7.5}{Pausar para o debate e ver a turma}{\pronto}
\hu{Professor}{parar a turma para debater e ver quais cláusulas foram
escolhidas}{discutir as escolhas sem transformá-las em acerto}
\aceite
\begin{itemize}
  \item \enquote{Pausar para Debate} (opcional) e o inverso.
  \item A distribuição das 4 cláusulas aparece em contagem, sem taxa de
    acerto.
  \item \enquote{Zerar a turma}.
\end{itemize}
\end{historia}

% ---------------------------------------------------------------------------
\feature{F06.8}{unooGPT --- Missão Tempestade}{\pronto}

\textbf{2ª série EM, Unidade 1} \textperiodcentered{} imita um assistente de
IA durante toda a missão \textperiodcentered{} uma aula \textperiodcentered{}
IA conversa com o estudante.

\textbf{Proposta.} Letramento de dados: perguntar melhor, perceber que fontes
discordam e decidir sob incerteza, sem confundir \enquote{não sabemos} com
\enquote{zero}. Nenhuma pergunta produz a resposta certa.

\begin{historia}{H06.8.1}{Dar a largada para a turma}{\pronto}
\hu{Professor}{controlar o relógio da turma}{que todos comecem e terminem
juntos}
\aceite
\begin{itemize}
  \item O cronômetro fica parado até \enquote{Dar a largada}.
  \item Também: +2 min, encerrar, modo livre e redefinir.
\end{itemize}
\end{historia}

\begin{historia}{H06.8.2}{Investigar perguntando}{\pronto}
\hu{Estudante do Médio}{fazer perguntas ao assistente sobre a tempestade que
chega}{decidir para onde mandar as equipes de resgate}
\aceite
\begin{itemize}
  \item Tenho 3 perguntas em 5 minutos para um \enquote{Cientista de Dados} que
    consulta 8 fontes, cada uma com um viés que eu não vejo.
  \item O mapa só acende o que as fontes consultadas cobrem; bairro em branco
    quer dizer \enquote{não sabemos}.
  \item Cada resposta leva de 10 a 30 segundos e usa o meu limite diário de
    IA.
\end{itemize}
\end{historia}

\begin{historia}{H06.8.3}{Decidir e ver o resultado como faixa}{\pronto}
\hu{Estudante do Médio}{mandar as equipes e ver o que aconteceu}{entender que
uma decisão sob incerteza tem vários desfechos possíveis}
\aceite
\begin{itemize}
  \item Duas equipes para quatro bairros; justifico com as fontes que abri ou
    com \enquote{foi intuição}.
  \item O resultado aparece como uma faixa de 20 repetições da mesma decisão,
    mais um fato concreto e a lista do que eu nunca abri. Sem pontuação.
  \item Posso recomeçar.
\end{itemize}
\end{historia}

\begin{historia}{H06.8.4}{Ler a turma e cada estudante}{\pronto}
\hu{Professor}{ver como cada estudante usou o tempo e as fontes}{discutir o
processo de decisão, não o placar}
\aceite
\begin{itemize}
  \item A folha da turma mostra a conduta, sem coluna de \enquote{resgatados}.
  \item O relatório de cada estudante mostra \enquote{Rodadas enfrentadas},
    \enquote{Até a primeira fala}, \enquote{Mudanças de destino},
    \enquote{O que a missão mediu}, \enquote{Como usou o tempo} e a conversa
    inteira.
  \item O guia traz a tabela dos 8 vieses, com o aviso de não projetá-la.
\end{itemize}
\pendencia{a ficha e o guia não passaram por teste em sala; estão marcados
como ponto de partida para revisão pedagógica.}
\end{historia}

% ---------------------------------------------------------------------------
\feature{F06.9}{Vila do Lume --- A Noite Longa}{\pronto}

\textbf{Para qualquer ano} (6 a 13 anos) \textperiodcentered{} jogo 3D próprio,
em primeira pessoa \textperiodcentered{} uma aula.

\textbf{Proposta.} Na camada visível, reconhecer um golpe que não tem os sinais
que se costuma ensinar a procurar. Na camada invisível, um retrato de valores
feito a partir de escolhas que custam algo --- e que define de que lado cada
estudante começa no Instanooria.

\begin{historia}{H06.9.1}{Acender as fogueiras da vila}{\pronto}
\hu{Estudante}{explorar a vila e acender as oito fogueiras}{salvar a vila antes
que a noite acabe}
\aceite
\begin{itemize}
  \item Uma instrução só: \enquote{A vila apagou. Acenda as oito fogueiras
    antes que a noite acabe.}
  \item Pelo caminho há gestos opcionais (ajudar um aldeão, acender um
    lampião, gravar o nome na pedra, tocar o sino), que não valem ponto.
  \item O nome da criança nunca é gravado.
\end{itemize}
\end{historia}

\begin{historia}{H06.9.2}{Receber mensagens no meio do jogo}{\pronto}
\hu{Estudante}{receber mensagens enquanto jogo}{descobrir, no debate, em quais
eu confiei}
\aceite
\begin{itemize}
  \item Chegam 5 mensagens sem pausar o jogo: verdadeiras e um golpe principal
    num de três níveis (grosseiro, autoridade com urgência, colega pedindo um
    passo pequeno).
  \item O nível é distribuído de forma equilibrada na turma; do 1º ao 5º ano o
    nível intermediário não é usado.
  \item O jogo não tem tela de resultado: a revelação é do professor.
\end{itemize}
\end{historia}

\begin{historia}{H06.9.3}{Ler a turma sem rotular ninguém}{\pronto}
\hu{Professor}{entender o que aconteceu com as mensagens na
turma}{conduzir o debate \enquote{Em quem dava pra confiar?}}
\aceite
\begin{itemize}
  \item Na revelação, projeto as mensagens das três versões lado a lado.
  \item Resultados: \enquote{As três bolhas nesta turma}, \enquote{O que
    aconteceu com as mensagens} e \enquote{Estudante por estudante} --- sem
    nenhum valor pessoal na lista.
  \item O relatório de uma criança abre pelo nível do golpe que ela recebeu e
    mostra \enquote{Os dez valores}; só há número quando há dois ou mais
    sinais, e \enquote{não teve a oportunidade} é diferente de \enquote{teve e
    não foi}.
  \item Partida abandonada não gera retrato, e a tela explica por quê.
  \item \enquote{Zerar a turma} (\enquote{Isso não tem volta.}) também
    redistribui as versões.
  \item A gestão vê o resumo agregado; a linha de uma criança, não.
\end{itemize}
\pendencia{o retrato de valores é uma hipótese de desenho, não um instrumento
validado. As reações dentro do jogo ainda não contam no limite diário de IA.}
\end{historia}

% ---------------------------------------------------------------------------
\feature{F06.10}{As missões que faltam}{\afazer}

A trilha tem 48 unidades (12 anos $\times$ 4). Estão publicadas 8, mais a Vila
do Lume. O 2º, o 3º e o 9º ano do Fundamental e a 1ª e a 3ª série do Médio
ainda não têm nenhuma; o 5º ano ficou sem a sua de 25 de setembro de 2026 até a
versão nova de Quem mandou isso? voltar (\id{F06.11}). As unidades
em produção aparecem para o professor como \enquote{em produção}, sem título
inventado. A próxima já decidida é
\textbf{A Seleção}, no 7º ano, Unidade 1, no lugar do Eco no Grupo.

\begin{historia}{H06.10.1}{Cada missão nova nasce completa}{\afazer}
\hu{PO}{que toda missão nova entre na trilha pelo mesmo caminho}{que a décima
missão custe menos que a nona, e nenhuma quebre uma promessa do produto}
\aceiteproposta
\begin{itemize}
  \item Tem ano e unidade, e aparece na trilha e na grade da sala.
  \item Só abre na turma em que o professor liberou.
  \item Aparece no painel da Unooria e \textbf{desliga} quando desligada
    (\id{H10.1.2}).
  \item Tem ficha e guia da aula com as nove seções, e a habilidade da BNCC.
  \item A correção é da plataforma; os resultados não têm ranking nem coluna de
    acerto.
  \item O texto livre da criança é lido só pelo professor da turma.
  \item Se usa IA: nada gerado chega à turma sem aprovação, há material pronto
    se a IA falhar, e o uso conta no limite diário.
  \item Toda a interface em português, exceto os termos do aplicativo imitado.
  \item Tem teste automatizado que percorre a missão como o estudante e como o
    professor.
\end{itemize}
\end{historia}

\begin{historia}{H06.10.2}{A Seleção (7º ano, Unidade 1)}{\afazer}
\hu{Estudante do 7º ano}{uma missão na rede social de fotos para a Unidade 1
do 7º ano}{que a trilha do 7º ano siga o material do 1º bimestre}
\aceiteproposta
\begin{itemize}
  \item Substitui por completo o Eco no Grupo; os dados do Eco são limpos antes.
  \item As imagens dos posts são geradas no preparo do professor, com
    alternativa pronta se a geração falhar.
  \item O detalhamento vem do documento de requisitos da missão.
\end{itemize}
\end{historia}

\begin{historia}{H06.10.3}{Quem mandou isso?, versão nova (5º ano, Unidade 1)}{\parcial}
\hu{Estudante do 5º ano}{a versão nova da missão do grupo de mensagens na
Unidade 1 do 5º ano}{aprender a conferir antes de repassar, agora também em
equipe}
\aceite
\begin{itemize}
  \item Substitui por completo a primeira versão, que saiu da trilha com os
    dados dela em 25 de setembro de 2026, e volta ao mesmo lugar da trilha.
  \item As mensagens do grupo se dividem em verdadeiras ou falsas e em
    emocionais ou neutras, e o estudante escolhe exatamente três para repassar.
  \item Depois de escolher, ele diz numa enquete o que pesou na escolha, e
    fecha com a regra \enquote{Se \ldots, então \ldots} e duas perguntas da
    Nina.
  \item A missão registra o que cada estudante demonstrou das habilidades
    (\id{E14}) e aceita equipes num computador dividido (\id{E13}).
  \item O detalhe, história por história, está em \id{F06.11}.
\end{itemize}
\pendencia{o teste que percorre a missão como o estudante e como o professor
foi escrito e ainda não rodou; a régua das habilidades aguarda a validação da
equipe pedagógica; e os textos escritos pela equipe aguardam o PO.}
\end{historia}
